% output=pdftex

\usemodule[res-05]

% medicijnen anchor ipv meer layers figbase : grid 

% drawgrid (boundingbox currentpicture, 10pt, 10pt, ticks) ...
% grid als overlay 
% grid=ja/nee,optie=test

\startMPinclusions 
  def drawgrid   expr e = do_drawgrid (e,false) enddef; 
  def drawulgrid expr e = do_drawgrid (e,true ) enddef; 
  def drawllgrid expr e = do_drawgrid (e,false) enddef; 
  def do_drawgrid (expr e, b) text t =
    begingroup ; save dx, dy, cp ;
    path cp ; cp := boundingbox currentpicture ;
    (dx,dy) = paired (e) ; 
    draw image (do_do_drawgrid(dx, dy, cp, b) t) ; 
    endgroup ; 
  enddef ;    
  def do_do_drawgrid (expr dx, dy, cp, b) text t = 
    if dy>0 : 
      def shft (expr i) = shifted (0,if b : bbheight(cp) - fi i) enddef ;
      for i := 0 step dy until bbheight(cp) : 
        draw (llcorner cp--lrcorner cp) shft(i) t ; 
      endfor ; 
      for i := 0 step 5dy until bbheight(cp) : 
        draw (llcorner cp shifted (-.50dy,0) -- 
              llcorner cp shifted (-.25dy,0)) shft(i) t ; 
        draw (lrcorner cp shifted ( .50dy,0) -- 
              lrcorner cp shifted ( .25dy,0)) shft(i) t ; 
      endfor ; 
    fi ; 
    if dx>0 : 
      def shft (expr i) = shifted (i,0) enddef ;
      for i := 0 step dx until bbwidth(cp) : 
        draw (llcorner cp--ulcorner cp) shft(i) t ; 
      endfor ; 
      for i := 0 step 5dx until bbheight(cp) : 
        draw (llcorner cp shifted (0,-.50dx) -- 
              llcorner cp shifted (0,-.25dx)) shft(i) t ; 
        draw (ulcorner cp shifted (0, .50dx) -- 
              ulcorner cp shifted (0, .25dx)) shft(i) t ; 
      endfor ; 
    fi ;  
    draw cp t ;
    setbounds currentpicture to cp ; 
  enddef ; 
\stopMPinclusions 

% \startuseMPgraphic{titlepage}
%   StartPage ;
%     path p, q, r ; pair a ;
%     fill Page withcolor .3green ; 
%     set_grid (PaperWidth, PaperHeight,50pt,50pt) ;
%     forever :
%       a := center Page randomized 2center Page ; 
%       if new_on_grid(xpart a,ypart a) :
%         a := (dx,dy) ; 
%         p := ((-1,-1) -- ( 1,1)) scaled 50pt shifted a ;
%         q := (( 1,-1) -- (-1,1)) scaled 50pt shifted a ;
%         r := origin shifted a ; 
%         pickup pencircle scaled 75pt ; 
%         draw r withcolor .3red ; 
%         pickup pencircle scaled 25pt ; 
%         draw p withcolor .3green ;
%         draw q withcolor .3green ;
%         pickup pencircle scaled 12.5pt ; 
%         draw p withcolor .4white ;
%         draw q withcolor .4white ;
%       fi ;
%       exitif grid_full ;
%     endfor ; 
%   StopPage ; 
% \stopuseMPgraphic

\startuseMPgraphic{titlepage}
  def SomeLabel (expr u) =
    path p ; 
    p := unitsquare xyscaled (10u,5u) smoothed (.5u) ;
    fill p withcolor c ; 
    draw p withcolor .5c withpen pencircle scaled .5u ; 
    drawoptions(withcolor .5c withpen pencircle scaled .25u) ; 
    for i=u step u until 3u : draw (2u,i)--(9u,i) ; endfor ;  
    p := fullcircle scaled u shifted (u,4u) ;
    drawoptions() ; 
    fill p withcolor .75c ;  
    draw p withpen pencircle scaled .25u withcolor .5c ;  
  enddef ; 
  StartPage ;
    picture pic ; pair loc ; color c ; c := (.8,.8,.45) ; 
    fill Page withcolor c ; 
    set_grid (PaperWidth, PaperHeight,4cm,4cm) ;
    forever :
      loc := center Page randomized 2center Page ; 
      if new_on_grid(xpart loc,ypart loc) :
        pic := image(SomeLabel(1cm)) ; 
        draw pic rotatedaround(center pic, -30+uniformdeviate 60)   
          shifted (dx,dy) shifted - center pic 
      fi ;
      exitif grid_full ;
    endfor ; 
  StopPage ; 
\stopuseMPgraphic

% @@figgd

\environment mstyle 

\starttext

\TitlePage{0}{0}{Adding Text\\to Graphics}{}{1}

\subject[introduction]{Introduction}

This is short manual about adding text to graphics made by
other applications than \TEX. Early versions of \CONTEXT\
already had provisions for adding information to graphics.
The \type {\start...\stopfigure} environment provides a way
to add text and hyperlinks to graphics based on a grid. We
used this feature to create interactive maps, navigate
using diagrams and alike. The corresponding definitions are
stored per graphic, and can be managed independently from
the main text. 

The method described in this document reimplements this
feature in a more flexible way using a couple of features
not present at that time. Also, the new method is more
suited to handle information stored in database like the
figure and resource databases supported by \CONTEXT. In due
time the old mechanism will be replaced by (i.e.\ redefined
in) the new one. 

The reason for extending the figure database concept with
this kind of information is that the people responsible for
the content (text) are not always the same as those making
the graphics. In some of our projects, authors are supposed
to add text (here called labels) to graphics. The same
graphic can be used in more than one context, with
different labels. Think for instance of a graphic that is
used in a question without labels, but in an answer with
labels. Or consider the same graphic being used in a Dutch
and English document. 

One handicap in separating graphic design and writing text
is that both the graphic designer and the author must make
sure that they know where the information ends up. Graphic
designers use professional drawing packages that authors
don't have access to, or demand in||depth knowlegde of the
application. Authors on the other hand know how to use
\TEX\ to typeset math, and drawing applications seldom
provide proper support for math. Separating drawing the
graphic and defining the labels also has the advantage that
the labels can be typeset in a way that suits the document
style (and specifically the fonts that are used). 

Although maintaining label specific data, like for instance
the locations where labels have to end up, is possible as
an independent activity, it may give the artist an uneasy
feeling, especially because he is used to click and point
tools. Therefore we will also discuss how to interface to 
Adobe Illustrator, a popular drawing application.  

\subject{Grids} 

Imagine that we have the following graphic defined in a 
drawing program. 

\startlinecorrection[blank]
\startMPcode
  loadfigure "labels-1.mp" number 1 ; 
\stopMPcode
\stoplinecorrection

If you want to add some texts to this graphic, you need to 
know where these should be anchored. One way to achieve this
is to put an imaginary grid on top of the graphic and anchor
labels at fixed positions. Because graphics can be included 
at different sizes, such a grid may change accordingly. 
Imagine that you would have to define labels using the grids
of the following graphics.

\startlinecorrection[blank]
\startcombination[3*1]
  {\startMPcode
     loadfigure "labels-1.mp" number 1 ; 
     drawgrid 10bp withpen pencircle scaled .5pt ;
   \stopMPcode}
  {natural size}
  {\startMPcode
     loadfigure "labels-1.mp" number 1 scaled .75 ; 
     drawgrid 10bp withpen pencircle scaled .5pt ;
   \stopMPcode}
  {scaled 75\%}
  {\startMPcode
     loadfigure "labels-1.mp" number 1 scaled 1.30 ; 
     drawgrid 10bp withpen pencircle scaled .5pt ;
   \stopMPcode}
  {scaled 130\%}
\stopcombination
\stoplinecorrection

In practice you will define points on a fixed grid layed 
over the graphic scaled at 100\%. In that case the grid will
scale with the graphic. 

\startlinecorrection[blank]
\startcombination[3*1]
  {\startMPcode
     loadfigure "labels-1.mp" number 1 ; 
     drawgrid 10bp withpen pencircle scaled .5pt ;
   \stopMPcode}
  {natural size}
  {\startMPcode
     loadfigure "labels-1.mp" number 1 ; 
     drawgrid 10bp withpen pencircle scaled (.5pt/.75) ;
     currentpicture := currentpicture scaled .75 ; 
   \stopMPcode}
  {scaled 75\%}
  {\startMPcode
     loadfigure "labels-1.mp" number 1 ; 
     drawgrid 10bp withpen pencircle scaled (.5pt/1.30) ;
     currentpicture := currentpicture scaled 1.30 ; 
   \stopMPcode}
  {scaled 130\%}
\stopcombination
\stoplinecorrection

Most drawing programs put their reference points in the
lower left corner. This makes sense since that suits
traditional coordinate systems. However, in a text flow it
makes more sense to think top||down. 

\startlinecorrection[blank]
\startcombination[2*1]
  {\startMPcode
     loadfigure "labels-1.mp" number 1 ; 
     drawulgrid 10bp withpen pencircle scaled .5pt ;
   \stopMPcode}
  {top anchored}
  {\startMPcode
     loadfigure "labels-1.mp" number 1 ; 
     drawllgrid 10bp withpen pencircle scaled .5pt ;
   \stopMPcode}
  {bottom anchored}
\stopcombination
\stoplinecorrection

The labeling mechanism described here works bottom up,
which is opposite to the default top||down placement in
\CONTEXT\ text layers. 

\subject{Adding text labels}

\resetfigurelabels[labels-1]
\resetfigurelabels[labels-2]

{\em This is a preliminary description. Multiple language
and label sets will be discussed as soon as we consider the
interface stable. We will also support other interfaces and 
ways of positioning.} 

If you have to figure out the positions on your own, the 
following method can be used to add labels to a graphic. 

\startbuffer[lab]
\startfigurelabels[labels-1]
  \definefigurelabel[x=25bp,y=45bp]{\bfd\white A}
  \definefigurelabel[x=70bp,y=30bp]{\bfd\white B}
  \definefigurelabel[x=60bp,y=75bp]{\bfd\white C}
\stopfigurelabels
\stopbuffer

\typebuffer[lab] 

The graphic itself is placed in the usual way: 

\startbuffer[fig]
\startlinecorrection[blank]
\externalfigure[labels-1.mps][option=label]
\stoplinecorrection
\stopbuffer

\typebuffer[fig] \getbuffer[lab] \getbuffer[fig]

or:  

\startbuffer[flt]
\placefigure
  {A floating figure}
  {\externalfigure[labels-1.mps][option=label,width=3cm]}
\stopbuffer

\typebuffer[flt] \getbuffer[lab] \getbuffer[flt] 

Although we limit ourselves here to simple labels, you can 
in principle put anything reasonable in a label. 

\subject{Using symbolic positions}

\resetfigurelabels[labels-1]
\resetfigurelabels[labels-2]

Especially when a graphic is used more than once with
different labels, or when the task of defining the anchors
can be delegated to the graphic designer, the separation
between defining anchors and texts comes into view. 

Each anchor gets a label (in its simplest form a number) and 
the positions are stored in a database. A record (which 
itself can be part of a figure (resource) library). 

[testen: pos in fig database] 

\startbuffer[fig]
\startlinecorrection[blank]
\externalfigure[labels-2.mps][option=label]
\stoplinecorrection
\stopbuffer

\typefile{labels-2.xml}

Here we have defined 4 positions that belong to figure \type
{labels-2}. 

\typebuffer[fig]  

\usefigurelabelbase[reset] \usefigurelabelbase[labels-2] 

{\showfigurelabelstrue\getbuffer[fig]}

It is possible to combine external and internal
definitions, so you can use an external \XML\ position database
and define the label texts in the document itself. The
label text definitions can be given in a \TEX\ syntax or in
\XML. 

The database can also contain the text labels themselves, 
like: 

\startbuffer
<rl:textlabels label="labels-2">
  <rl:textlabel label="1">A</rl:textlabel>
  <rl:textlabel label="2">B</rl:textlabel>
  <rl:textlabel label="3">C</rl:textlabel>
  <rl:textlabel label="4">D</rl:textlabel>
</rl:textlabels>
\stopbuffer

\typebuffer \processXMLbuffer {\getbuffer[fig]}

\resetfigurelabels[labels-2]

If the source document is a normal \TEX\ document, you can 
include the definitions in your file. 

\startbuffer
\startfigurelabels[labels-2]
  \definefigurelabel[1]{\bfd\white A}
  \definefigurelabel[2]{\bfd\white B}
  \definefigurelabel[3]{\bfd\white C}
  \definefigurelabel[4]{\bfd\white D}
\stopfigurelabels
\stopbuffer

\typebuffer \getbuffer {\getbuffer[fig]}

You can add additional labels. If needed you can provide 
your own co\"ordinates. 

\startbuffer
\startfigurelabels[labels-2]
  \definefigurelabel[x=50bp,y=50bp]{\bfd\white E}
  \definefigurelabel[1]{\bfd\red\symbol[star]}
\stopfigurelabels
\stopbuffer

\typebuffer \getbuffer {\getbuffer[fig]}

If you want a fresh start, you should explicitly reset the 
data with the reset command. We use this options to show you
the alternative alignment locations (these are the same as 
in the \CONTEXT\ layer mechanism). 

\startbuffer
\resetfigurelabels[labels-2]
\startfigurelabels[labels-2]
  \definefigurelabel[1][location=l]{l}
  \definefigurelabel[2][location=r]{r}
  \definefigurelabel[3][location=t]{t}
  \definefigurelabel[4][location=b]{b}
\stopfigurelabels
\stopbuffer

\typebuffer \getbuffer {\traceboxplacementtrue \getbuffer[fig]}

\startbuffer
\resetfigurelabels[labels-2]
\startfigurelabels[labels-2]
  \definefigurelabel[1][location=lt]{lt}
  \definefigurelabel[2][location=lb]{lb}
  \definefigurelabel[3][location=rt]{rt}
  \definefigurelabel[4][location=rb]{rb}
\stopfigurelabels
\stopbuffer

\typebuffer \getbuffer {\traceboxplacementtrue \getbuffer[fig]}

\subject{Adobe Illustrator files}

\resetfigurelabels[labels-1]
\resetfigurelabels[labels-2]

The WARM plugin of Adobe Illustrator gives you the means to
tag positions in a graphic. This plug||in is an initiative 
by Ross Moore and Wendy Mackay. The WARM plug||in writes
special comment||only \POSTSCRIPT\ files. This means that
we don't have to ask graphic artists to use text base tools
for providing the positional information. 

The positions are normally numbered, but you can give them
meaningful names. These positions are saved in files with
the suffix \type {bb}. In Illustrator, these positions are
named marked points. The data segment of such a file looks
as follows: 

\typefile{labels-1.bb}

For our purpose, Only the coordinate (first entry) and
label (third entry) make sense. There can best be some
logic in placing the points, especially since we have to
align the labels manually. When applied to our graphic,
the previous definitions result in the following label 
anchors. 

\startbuffer[pair]
\startcombination[3*1]
  {\externalfigure[labels-1.mps][option=label]}            {natural size}
  {\externalfigure[labels-1.mps][option=label,scale=750]}  {scaled 75\%}
  {\externalfigure[labels-1.mps][option=label,scale=1250]} {scaled 125\%}
\stopcombination
\stopbuffer

\startbaselinecorrection
  \showfigurelabelstrue \getbuffer[pair]
\stopbaselinecorrection

As you can see, the positions are scaled with the graphic, 
but the same \type {bb} is used for each of them. 

\startbaselinecorrection
  \showfigurelabelstrue \tracefigurelabelstrue \getbuffer[pair]
\stopbaselinecorrection

We define some labels:

\startbuffer
\startfigurelabels[labels-1]
  \definefigurelabel[1][location=r]{kwik}
  \definefigurelabel[2][location=r]{kwek}
  \definefigurelabel[3][location=l]{kwak}
\stopfigurelabels
\stopbuffer

\typebuffer \getbuffer 

These show up as follows. Watch how we aligned them left and
right of the anchor point. 

\startbaselinecorrection
  \getbuffer[pair]
\stopbaselinecorrection

Internally, \TEX\ works with real points, like \type {12pt},
but if needed you can define positions in \POSTSCRIPT\ points, like 
\type {12bp}. Watch out: these are not the same, although 
for applications like these the difference does not show of 
that fast. 

\startbaselinecorrection
  \centerline
    {\tracefigurelabelstrue
     \externalfigure[labels-1.mps][option=label]}      
\stopbaselinecorrection

It is quite possible that the author wants to put a couple
of extra labels in a graphic. 

\startbuffer
\startfigurelabels[labels-1]
  \definefigurelabel[x=50bp,y=50bp]{kwok}
\stopfigurelabels
\stopbuffer

\typebuffer \getbuffer

The entries are added to the already defined ones; if you 
want to start fresh, you should explicitly reset the label 
texts with: 

\starttyping
\resetfigurelabels[labels-1]
\stoptyping

\startbaselinecorrection
  \getbuffer[pair]
\stopbaselinecorrection

The WARM method is hooked into the external figure mechanism
as (optional) second step in resolving layers. This means 
that an existing \XML\ definition (file) takes precedence. 
In any case, the way to invoke this feature is the same: 

\starttyping 
\externalfigure[name][option=label]
\stoptyping 

In combination with previously defined labels this will give
you labeled figures, given that a \type {bb} file is 
present. You can convert such files to an \XML\ file using the 
\type {bbtoxml} \PERL\ script. The generated base can be 
registered by saying: 

\starttyping 
\usefigurelabelbase[reset] \usefigurelabelbase[bbtoxml]
\stoptyping 

\ColofonPage

\stoptext 

% labels-1.mp 

input mp-tool ; 

beginfig (1) ; 
  path p, q[] ; color gray ; 
  red   := 1.4(.6,.4,.5) ; 
  green := 1.4(.4,.6,.5) ; 
  blue  := 1.4(.4,.5,.6) ; 
  gray  := 1.4(.5,.5,.5) ; 
  pickup pencircle scaled 5pt ; 
  p := fullcircle scaled 100pt ; 
  q[1] := center p -- point 0 of p --
          (p cutbefore point 0 of p) cutafter point 3 of p --
          point 3 of p -- cycle ; 
  q[2] := center p -- point 3 of p --
          (p cutbefore point 3 of p) cutafter point 6 of p --
          point 6 of p -- cycle ; 
  q[3] := center p -- point 6 of p --
          (p cutbefore point 6 of p) cutafter point length(p) of p -- 
          point length(p) of p -- cycle ; 
  fill q[1] withcolor red ; draw q[1] withcolor gray ;
  fill q[2] withcolor green ; draw q[2] withcolor gray ;
  fill q[3] withcolor blue ; draw q[3] withcolor gray ;
endfig ; 

end .

% labels-2.mp

input mp-tool ; 

beginfig (1) ; 
  path p, q[] ; color gray ; 
  red   := 1.4(.6,.4,.5) ; 
  green := 1.4(.4,.6,.5) ; 
  blue  := 1.4(.4,.5,.6) ; 
  gray  := 1.4(.5,.5,.5) ; 
  pickup pencircle scaled 5pt ; 
  p := fullcircle scaled 100pt ; 
  q[1] := center p -- point 0 of p --
          (p cutbefore point 0 of p) cutafter point 3 of p --
          point 3 of p -- cycle ; 
  q[2] := center p -- point 3 of p --
          (p cutbefore point 3 of p) cutafter point 6 of p --
          point 6 of p -- cycle ; 
  q[3] := center p -- point 6 of p --
          (p cutbefore point 6 of p) cutafter point length(p) of p -- 
          point length(p) of p -- cycle ; 
  fill q[1] withcolor red ; draw q[1] withcolor gray ;
  fill q[2] withcolor green ; draw q[2] withcolor gray ;
  fill q[3] withcolor blue ; draw q[3] withcolor gray ;
  label ("x", origin) ; 
endfig ; 

end .
